\documentclass[12pt]{article}
\usepackage[utf8]{inputenc}
\usepackage[T1]{fontenc}
\usepackage[margin=1in]{geometry}
\usepackage{setspace}
\usepackage{xcolor}
\usepackage[colorlinks=true, urlcolor=blue!60!black, linkcolor=blue!60!black]{hyperref}
\setlength{\parindent}{0em}
\setlength{\parskip}{0.45em}
\renewcommand{\baselinestretch}{1.0}
\begin{document}
\begin{center}
{\Large \textit{Martin v. Hunter's Lessee}, 14 U.S. (1 Wheat.) 304 (1816)\par}
\vspace{0.3em}
{\normalsize Case Brief --- One-Page Edition\par}
\end{center}

\textbf{Memory Jogger.} The Supreme Court may review state-court judgments on federal questions (\href{https://supreme.justia.com/cases/federal/us/14/304/}{Justia}).

\textbf{Facts.} Lord Fairfax devised Virginia lands to his British nephew; Virginia granted a parcel to David Hunter. Federal treaties of 1783 and 1794 protected British subjects' lands. Chief Justice Marshall, financially interested, recused.

\textbf{Procedural History.} After the Supreme Court upheld Fairfax's title in 1813, Virginia's Court of Appeals refused to obey the mandate, holding Section~25 of the Judiciary Act unconstitutional as applied to state courts.

\textbf{Votes.} Reversed, 6--0 (Story for the Court; Johnson concurring; Marshall recused).

\textbf{Holding.} Article~III appellate jurisdiction reaches state-court judgments on federal questions because it attaches to the case, not the court; Section~25 is valid, and a refusal to obey a mandate is itself reviewable.

\textbf{Reasoning.} The people, not the States, ordained the Constitution, and Article~III's appellate grant has no limit on the tribunal appealed from. The Supremacy Clause contemplates state judges deciding federal questions, local bias may sway them, and uniformity demands one revising court; the First Congress agreed. \emph{Separately:} Johnson stressed that the Court acts ``over persons and cases,'' not over state courts as subordinates.

\textbf{Examples.} \emph{Same-Side:} a state conviction rejecting a federal defense. \emph{Opposite-Side:} a judgment on an adequate and independent state ground. \emph{Fence-Sitter:} congressional stripping of all federal jurisdiction over a class of claims.

\textbf{Critique.} Attacked as ``commandeered'' state courts; celebrated as the linchpin of national rights.

\textbf{Key Quotation.} ``That motive is the importance, and even necessity, of uniformity of decisions throughout the whole United States upon all subjects within the purview of the Constitution. Judges of equal learning and integrity in different States might differently interpret a statute or a treaty of the United States, or even the Constitution itself.'' --- Justice Story, 14 U.S. at 347--48.

\end{document}
